\documentclass[12pt,a4paper]{article}

% ------------------------------------------------
% Packages
% ------------------------------------------------
\usepackage[utf8]{inputenc}
\usepackage[T1]{fontenc}

\usepackage{amsmath}
\usepackage{amssymb}

\usepackage{enumitem}
\usepackage{hyperref}

% ------------------------------------------------
% Bibliography
% ------------------------------------------------
\usepackage[
    backend=biber,
    style=numeric,
    sorting=none
]{biblatex}

\addbibresource{reference.bib}

% ------------------------------------------------
% Title
% ------------------------------------------------
\title{
    Retrieval Quality as the Primary Failure Point in RAG-Based and Agentic Research Assistants
}

\author{Abhishek}
\date{}

% ------------------------------------------------
% Document
% ------------------------------------------------
\begin{document}

\maketitle

% =================================================
\section{Abstract}
% =================================================

Retrieval-Augmented Generation (RAG) systems combine information retrieval
with large language models to improve the factual grounding of generated
responses. However, retrieval quality remains a critical determinant of
overall system reliability. When relevant, sufficient, authoritative, and
temporally valid evidence is not retrieved, even a capable language model
may generate an incomplete or incorrect answer.

Recent developments in agentic research assistants have further increased
the importance of retrieval because these systems perform multi-step
information-seeking, query reformulation, browsing, evidence synthesis,
and iterative reasoning. This paper examines retrieval quality as a primary
failure point in RAG-based and agentic research assistants. It reviews
retrieval failures, query formulation, evidence coverage, reranking,
verification, adaptive retrieval, and evidence-aware generation.

The paper argues that retrieval should not be treated merely as a
preprocessing stage but as an adaptive and evidence-aware component of the
complete research-assistant pipeline.

% =================================================
\section{Keywords}
% =================================================

Retrieval-Augmented Generation, RAG, Agentic RAG, Information Retrieval,
Evidence Retrieval, Query Reformulation, Retrieval Quality, Evidence
Verification, Research Assistants, Large Language Models

% =================================================
\section{Introduction}
\label{sec:introduction}
% =================================================

Retrieval-Augmented Generation (RAG) has emerged as an important approach
for improving the factual grounding and knowledge access capabilities of
large language models (LLMs). Instead of relying exclusively on knowledge
encoded in model parameters, RAG systems retrieve external information and
provide the retrieved evidence to the generation model
\cite{lewis2020rag}.

The retrieval component therefore plays a fundamental role in determining
the quality of the final response. A generation model cannot reliably use
information that was never retrieved. Consequently, retrieval failures can
propagate directly into generation failures.

A useful approximation of this dependency can be expressed as shown in
Equation~\ref{eq:reliability-decomposition}.

\begin{equation}
P(\text{correct answer})
\approx
P(\text{retrieve sufficient evidence})
\times
P(\text{use evidence correctly} \mid \text{evidence retrieved})
\times
P(\text{generate correctly} \mid \text{evidence used}).
\label{eq:reliability-decomposition}
\end{equation}

This formulation highlights retrieval as an important reliability
bottleneck. Even when the language model has strong reasoning and
generation capabilities, insufficient evidence retrieval can constrain the
quality of the final answer.

Traditional dense retrieval methods have substantially improved semantic
search. Dense Passage Retrieval (DPR), for example, demonstrated the
effectiveness of dense representations for open-domain question answering
\cite{karpukhin2020dpr}. Evaluation benchmarks such as BEIR subsequently
highlighted the difficulty of generalizing retrieval systems across diverse
datasets and domains \cite{thakur2021beir}.

Recent research has also demonstrated that retrieving information is not
sufficient by itself. The position and distribution of retrieved evidence
can affect the ability of language models to use that evidence effectively
\cite{liu2024lostmiddle}. Query expansion and alternative query generation
methods have therefore become increasingly important for improving
retrieval coverage \cite{gao2023hyde,wang2023query2doc}.

The emergence of agentic research assistants further complicates the
retrieval problem. Agentic systems may execute multiple searches, interact
with web environments, reformulate queries, and iteratively collect
evidence. Systems such as WebArena and WebVoyager illustrate the increasing
importance of autonomous web interaction and information acquisition
\cite{zhou2024webarena,he2024webvoyager}.

This paper investigates retrieval quality as a primary failure point in
RAG-based and agentic research assistants and examines how adaptive,
evidence-aware retrieval mechanisms can improve system reliability. The
proposed reliability framework is developed in
Section~\ref{sec:framework} and the corresponding pipeline is presented in
Section~\ref{sec:pipeline}.

% =================================================
\section{Background and Related Work}
\label{sec:background}
% =================================================

\subsection{Retrieval-Augmented Generation}

RAG combines a retrieval mechanism with a generative language model.
The retrieval component searches an external knowledge source, while the
generator uses the retrieved information to produce an answer
\cite{lewis2020rag}.

The basic RAG pipeline can be represented as:

\[
\text{Query}
\rightarrow
\text{Retriever}
\rightarrow
\text{Retrieved Evidence}
\rightarrow
\text{Generator}
\rightarrow
\text{Answer}.
\]

Although this architecture is conceptually simple, retrieval quality
strongly affects the final answer.

\subsection{Dense Retrieval}

Dense retrieval represents queries and documents in vector space and
retrieves documents based on semantic similarity. DPR demonstrated that
dense passage retrieval can outperform traditional retrieval approaches
for several open-domain question-answering settings
\cite{karpukhin2020dpr}.

However, semantic similarity does not necessarily guarantee that retrieved
documents contain all information required to answer a question.

\subsection{Retrieval Evaluation}

The BEIR benchmark demonstrated that retrieval systems can behave
differently across datasets and domains, highlighting the importance of
robust and generalizable retrieval evaluation
\cite{thakur2021beir}.

Traditional retrieval metrics such as recall, precision, and ranking
quality are useful, but RAG systems require additional evaluation of
evidence sufficiency and downstream answer quality.

RAGAS introduced an evaluation framework for assessing RAG pipelines using
metrics related to faithfulness, answer relevance, and context relevance
\cite{es2024ragas}.

ARES similarly focuses on automated evaluation of retrieval-augmented
generation systems \cite{saadfalcon2024ares}.

% =================================================
\section{Retrieval as an Evidence Bottleneck}
\label{sec:bottleneck}
% =================================================

A central problem in RAG systems is that retrieval errors can propagate
through the remainder of the pipeline.

Let the required evidence for answering a question be:

\begin{equation}
E^{*} = \{e_1,e_2,\ldots,e_n\}.
\label{eq:required-evidence}
\end{equation}

Let the retrieved evidence be:

\begin{equation}
R = \{r_1,r_2,\ldots,r_k\}.
\label{eq:retrieved-evidence}
\end{equation}

For reliable answering, the retrieved information should provide adequate
coverage of the required evidence, as expressed in
Equation~\ref{eq:coverage}.

\begin{equation}
E^{*} \subseteq coverage(R).
\label{eq:coverage}
\end{equation}

When important evidence is absent, the generation model may be forced to
infer missing information, rely on parametric knowledge, or produce an
unsupported answer.

This problem becomes more significant for multi-hop questions where the
answer depends on multiple pieces of evidence.

% =================================================
\section{Retrieval Failure Modes}
\label{sec:failure-modes}
% =================================================

\subsection{Low Recall}

Low recall occurs when relevant information exists in the knowledge source
but is not retrieved. This is particularly problematic when the missing
evidence contains a critical fact required for answering the query.

\subsection{Low Precision}

Retrieval systems can also return large quantities of irrelevant evidence.
This increases the context size and can make it harder for the language
model to identify useful information.

\subsection{Incomplete Evidence Coverage}

A retrieval system may retrieve individually relevant documents without
retrieving the complete evidence chain required for a multi-hop answer.

\subsection{Evidence Position}

The location of relevant information within a long context can influence
model performance. The ``lost in the middle'' phenomenon demonstrates that
models may underutilize information placed in certain positions within
long contexts \cite{liu2024lostmiddle}.

\subsection{Temporal Staleness}

Information may become outdated over time. Therefore, retrieval quality
should consider the temporal validity of retrieved evidence.

\subsection{Source Authority}

Relevant information is not necessarily reliable information. Research
assistants therefore need mechanisms for assessing the authority and
credibility of sources.

% =================================================
\section{Query Formulation and Multi-Representation Retrieval}
\label{sec:query-formulation}
% =================================================

The quality of retrieval depends partly on the formulation of the search
query.

For a complex information need $Q$, the system may generate multiple
representations, as shown in Equation~\ref{eq:query-representations}.

\begin{equation}
Q \rightarrow \{q_1,q_2,\ldots,q_m\}.
\label{eq:query-representations}
\end{equation}

Each representation can target a different aspect of the information need.

HyDE demonstrated the use of hypothetical documents to improve dense
retrieval by transforming queries into representations that more closely
resemble retrieved documents \cite{gao2023hyde}.

Query2doc similarly explores query expansion using generated textual
representations \cite{wang2023query2doc}.

These approaches suggest that retrieval should not necessarily rely on a
single literal query representation.

% =================================================
\section{Current Approaches for Improving Retrieval}
\label{sec:current-approaches}
% =================================================

\subsection{Corrective Retrieval}

Corrective RAG introduces mechanisms for identifying poor retrieval results
and performing corrective actions when retrieved information is inadequate
\cite{yan2024crag}.

The central idea is that retrieval should be evaluated before generation
rather than blindly passed to the language model.

\subsection{Self-Reflective Retrieval}

Self-RAG introduces reflection mechanisms that allow a language model to
evaluate retrieved information and its own generation process
\cite{asai2024selfrag}.

This provides a mechanism for incorporating retrieval and self-reflection
into the generation process.

\subsection{Retrieval-Aware Ranking}

RankRAG investigates ranking and instruction-based approaches for improving
the use of retrieved information \cite{yu2024rankrag}.

Ranking is important because a retrieval system may return more evidence
than can be effectively processed by the generator.

\subsection{Retrieval Evaluation and Error Analysis}

RAGChecker provides a framework for fine-grained evaluation and diagnosis
of RAG systems \cite{ru2024ragchecker}.

RAGTruth focuses on detecting hallucination-related issues in
retrieval-augmented generation systems \cite{niu2024ragtruth}.

These approaches demonstrate the importance of evaluating retrieval and
generation as interconnected components.

% =================================================
\section{Agentic Research Assistants}
\label{sec:agentic}
% =================================================

Agentic research assistants extend conventional RAG systems by introducing
planning, iterative retrieval, web browsing, tool use, and multi-step
reasoning.

Systems designed for complex research tasks may perform several retrieval
operations rather than executing a single search.

STORM demonstrates an approach for automated research and article
generation through structured information gathering and multi-perspective
search \cite{shao2024storm}.

Recent research has also explored search-oriented reasoning and deep
research systems \cite{li2025searcho1,zheng2025deepresearcher,
schmidgall2025agentlaboratory}.

Web-based agent benchmarks such as WebArena and WebVoyager demonstrate the
complexity of autonomous interaction with web environments
\cite{zhou2024webarena,he2024webvoyager}.

Mind2Web further demonstrates the challenge of using language models for
web interaction and task execution \cite{deng2023mind2web}.

These systems reinforce the importance of reliable retrieval across
multiple steps of an agentic workflow, as discussed further in
Section~\ref{sec:challenges}.

% =================================================
\section{Retrieval Reliability Framework}
\label{sec:framework}
% =================================================

A retrieval reliability framework can be represented as shown in
Equation~\ref{eq:reliability-tuple}.

\begin{equation}
E = (R,P,C,A,T,D,V,U,G),
\label{eq:reliability-tuple}
\end{equation}

where:

\begin{itemize}
    \item $R$ = retrieval recall,
    \item $P$ = retrieval precision,
    \item $C$ = evidence coverage,
    \item $A$ = source authority,
    \item $T$ = temporal validity,
    \item $D$ = source diversity,
    \item $V$ = evidence verification,
    \item $U$ = evidence utilization,
    \item $G$ = generation quality.
\end{itemize}

This formulation treats retrieval reliability as a multidimensional
property rather than a single retrieval metric.

% =================================================
\section{Proposed Retrieval Reliability Pipeline}
\label{sec:pipeline}
% =================================================

A reliable retrieval pipeline, building on the framework in
Section~\ref{sec:framework}, can be organized into the following stages.

\subsection{Information-Need Analysis}

The system first identifies the actual information requirement behind the
user's question.

\subsection{Multi-Representation Querying}

The original query is transformed into multiple query representations to
increase retrieval coverage.

\subsection{High-Recall Retrieval}

The retrieval stage prioritizes obtaining potentially relevant evidence
before aggressive filtering.

\subsection{Evidence-Aware Reranking}

Retrieved documents are reranked according to relevance, evidence
coverage, authority, and other reliability criteria.

\subsection{Evidence Verification}

The system checks whether the retrieved evidence actually supports the
information requirement.

\subsection{Adaptive Retrieval}

If the evidence is insufficient, the system performs additional retrieval,
query reformulation, or source exploration.

\subsection{Evidence-Constrained Generation}

The final response should be generated using verified evidence rather than
unsupported assumptions.

% =================================================
\section{Challenges and Limitations}
\label{sec:challenges}
% =================================================

Several challenges remain.

First, retrieval quality is difficult to evaluate using a single metric.
High recall does not necessarily imply that the retrieved evidence is
sufficient for answering a complex question.

Second, multi-hop information needs require evidence from multiple sources.
A retrieval system may retrieve relevant individual documents while still
failing to establish the complete evidence chain.

Third, source authority and temporal validity are difficult to measure
automatically.

Fourth, increasing retrieval depth may improve coverage but also increase
context size, computational cost, and irrelevant information.

Finally, agentic systems introduce additional failure opportunities because
errors can accumulate across multiple retrieval and reasoning steps.

% =================================================
\section{Research Gaps and Future Directions}
\label{sec:future-directions}
% =================================================

Future research should investigate adaptive retrieval systems that
dynamically determine when additional evidence is required.

An important direction is evidence-aware stopping. Instead of retrieving a
fixed number of documents, the system should determine whether the current
evidence set is sufficient for the information need.

Another important direction is source-aware retrieval. Retrieval systems
should incorporate source authority, temporal validity, diversity, and
cross-source agreement.

Future research should also investigate retrieval mechanisms specifically
designed for multi-hop and agentic research tasks.

Evaluation should move beyond conventional retrieval metrics toward
end-to-end measures of evidence sufficiency, verification, utilization,
and answer reliability.

% =================================================
\section{Conclusion}
\label{sec:conclusion}
% =================================================

Retrieval quality is a critical determinant of reliability in
retrieval-augmented and agentic research assistants. Retrieval failures can
propagate into downstream reasoning and generation, making it difficult for
even capable language models to produce reliable answers.

Existing research has introduced dense retrieval, query expansion,
corrective retrieval, self-reflection, ranking, retrieval evaluation, and
agentic web interaction. However, these approaches do not completely solve
the problem of determining whether the retrieved evidence is sufficient,
reliable, diverse, current, and properly utilized.

A more reliable research assistant therefore requires an adaptive,
evidence-aware retrieval pipeline that integrates information-need
analysis, multi-representation querying, high-recall retrieval,
evidence-aware reranking, verification, adaptive retrieval, and
evidence-constrained generation, as outlined in
Section~\ref{sec:pipeline}.

% =================================================
% References
% =================================================

\printbibliography

\end{document}
